\documentclass[10pt,a4paper]{amsart}
\usepackage[margin=19mm]{geometry}
\usepackage{amsmath,amssymb,mathtools}
\usepackage[scr=rsfs,cal=euler]{mathalfa}
\usepackage{xfrac}
\usepackage[unicode,hidelinks,bookmarks=false]{hyperref}
\newcommand{\ie}{i.e.\ }
\newcommand{\cf}{cf.\ }
\newcommand{\resp}{respectively\ }
\newcommand{\Subsection}{\subsection}
\providecommand{\nobreakdash}{\nobreak\hbox{-}}
\setlength{\emergencystretch}{2em}
\multlinegap0pt
\usepackage{bm}
\usepackage{bbm}
\usepackage{dsfont}
\usepackage{xspace}
\usepackage{cancel}
\usepackage{mathtools}
\usepackage{amsthm}
\makeatletter
\@ifundefined{thm}{\newtheorem{thm}{Thm}}{}
\@ifundefined{cor}{\newtheorem{cor}[thm]{Corollary}}{}
\@ifundefined{defn}{\newtheorem{defn}[thm]{Definition}}{}
\@ifundefined{lem}{\newtheorem{lem}[thm]{Lemma}}{}
\@ifundefined{prop}{\newtheorem{prop}[thm]{Proposition}}{}
\makeatother
\newcommand{\N}{\mathbb{N}}
\newcommand{\Q}{\mathbb{Q}}
\newcommand{\Z}{\mathbb{Z}}
\renewcommand{\P}{\mathbb{P}}
\title[Families of isogenous elliptic curves ordered by height]{Families of isogenous elliptic curves ordered by height}
\author{Yu Fu}
\date{Source: arXiv:2308.11122v3 (CC BY 4.0)}
\begin{document}
\maketitle

\noindent\emph{Source and licence.} Families of isogenous elliptic curves ordered by height. Version arXiv:2308.11122v3 (CC BY 4.0), distributed under the Creative Commons Attribution 4.0 International licence, as recorded in the arXiv licence metadata for this exact version and shown on the abstract page https://arxiv.org/abs/2308.11122v3. Full rights notice: licenses/arxiv-2308.11122-CC-BY-4.0.txt.\par\smallskip

\noindent\textit{Editorial scope.} Counted unit: the Prop whose source label is \texttt{change of heights} in the article (source0.tex lines 775--790, source label \texttt{change of heights}), delivered together with the article's own complete original proof of it (source0.tex lines 792--831, one \texttt{proof} environment of 40 lines). No mathematics is written, repaired or re-derived: statement and proof are the article's own text line for line. Required context retained uncounted: iota (lines 164--166); ProductHeights (lines 267--272); Hprod (lines 434--436); covers (lines 447--450); F (lines 571--576); iotaHp (lines 705--707); pazc (lines 742--746); 5.3 (lines 751--754). Every definition, display and label of those lines is retained unchanged. Omitted from this package and not redistributed: 1--163, 167--266, 273--433, 437--446, 451--570, 577--704, 708--741, 747--750, 755--774, 791--791, 832--1028. Nothing inside a retained excerpt is omitted or altered: every retained body is delivered exactly as the article's lines are written, so no omission receipt of any kind accompanies this package. Citations: the retained excerpts carry no citation command at all, so no bibliography entry of the article is transcribed and no reference is redistributed. Reference adaptation: none was needed and none was made; every reference inside the delivered unit resolves inside it, so no unresolved reference or citation is produced. Printed slips kept literal: no printed word, symbol, hypothesis or number of the counted statement or of its proof was altered. Layout and class: the article's own class is \texttt{amsart} with packages including amsmath, amsthm, amssymb, amsfonts, mathrsfs, amscd, tikz-cd, hyperref, ulem, xcolor, braket, fullpage, graphicx, multirow; the wrapper is the lane's pinned \texttt{amsart} editorial wrapper at 10pt on A4, which restates the theorem-like environments the retained text needs and adds the article's own macro definitions that the retained text needs (\texttt{\textbackslash N}, \texttt{\textbackslash P}, \texttt{\textbackslash Q}, \texttt{\textbackslash Z}), with extra wrapper packages (bm, bbm, dsfont, xspace, cancel, mathtools). The delivered numbering is the unit's own. Assets: no retained span contains a figure environment, a graphics-inclusion command, a PDF-page inclusion, a table or a TikZ picture; no image file is copied into this package, the package declares no asset dependency and no image file is redistributed with it. Licence: version arXiv:2308.11122v3 of this article is distributed under the Creative Commons Attribution 4.0 International licence, as recorded in the arXiv licence metadata for this exact version and shown on the abstract page https://arxiv.org/abs/2308.11122v3. A full rights notice is delivered at licenses/arxiv-2308.11122-CC-BY-4.0.txt. Characters: every retained excerpt is pure ASCII, so the delivered engine, XeLaTeX with its default Latin Modern text font, reports no missing character.
\begin{equation}\label{iota}
	\iota: C \rightarrow X(1) \times X(1)  \rightarrow \mathbb{P}^{1} \times \mathbb{P}^{1} \hookrightarrow \mathbb{P}^{3}.
	\end{equation}

\begin{lem}\label{ProductHeights}
Let $\sigma_{n}$ be the Segre embedding of $n$-copies of $\P^{1}$
$$\sigma_{n}\colon \underbrace{\P^{1} \times \P^{1} \times \ldots \times \P^{1}}_{n\text{-times}} \hookrightarrow \P^{2^{n}-1}.  $$  Let $H(.)$ be the projective height defined above. We have 
$$H(\sigma_{n}(x_1, \cdots, x_{n}))=H(x_{1})\cdots H(x_{n}).$$
	
\end{lem}

\begin{defn}\label{Hprod}
  Let $H_{\text{p}}:=H_{1} \times H_{2}$, where $H_{1}$ and $H_{2}$ (possibly $H_{1}=H_{2}$) are maximal parabolic subgroups of $SL_{2}(\Z/m\Z)$.
\end{defn}

\begin{lem}\label{covers}
	
	Any rational point $t \in S(B)$ for some $B$, where the specializations $E_t$ and $E^{\prime}_t$ above $t$ are non-CM, admits a lifting to one of the congruence covers: $X_{\tilde{H}_{\Delta}}(K)$ or $X_{H_{\text{p}}}(K)$.
\end{lem}

\begin{defn}\label{F}

	Let $F$ be a function defined on $X_{\tilde{H}_{\Delta}}$ given by
	$$F(E_{t}, E_{t}^{\prime}, \psi)=j(E_{t,1})j(E^{\prime}_{t,1})+ \cdots +j(E_{t, m+1})j(E^{\prime}_{t, m+1}).$$
	
	\end{defn}

\begin{equation}\label{iotaHp}
	\iota_{H_{\operatorname{p}}}: X_{0}(m) \times X_{0}(m) \hookrightarrow \P^{1} \times \P^{1} \times \P^{1} \times \P^{1} \hookrightarrow \P^{15}.
\end{equation}

\begin{cor}\label{pazc}
Let $\varphi: E_{1} \rightarrow E_{2}$ be a $\overline{\mathbb{Q}}$-isogeny between two elliptic curves defined over $\overline{\Q}$ which is cyclic of degree $m$. Let $j_{1}$ and $j_{2}$ be the respective $j$-invariants. Then one has
	$$H(j_{1}) < Am^{12}H(j_{2})$$
	for some absolute constant $A$. Here $H(.)$ denotes the projective height.
\end{cor}

\begin{lem}\label{5.3}
Let $d$ be the projective degree of $C$ under the embedding $\iota$, see (\ref{iota}). Let $H(\iota)$ be the height of $\iota$ defined by the height of the coefficients of the defining polynomials of $j(E_{t})$ and $j(E_{t}^{\prime}).$ Then for every $t \in K$ with $H(t) \le B$, we have
$$H(P_{t}) \le (d+1)H(\iota)B^{d}.$$	
\end{lem}

\begin{prop}\label{change of heights}

Fix $m \in \N$. Let $t \in K$ be a rational point such that $t \in S(B)$ for some $B$ and the specializations $E_t$ and $E^{\prime}_t$ above $t$ are non-CM elliptic curves. Let $P_{t}$ denote the point on $C$ parametrized by $t$, and denote by $\tilde{P_{t}}$ a lifting of $P_{t}$ to one of the covers $C_H$ in Lemma \ref{covers}. Let $H(\iota_{H}(\tilde{P_{t}}))$ denote the projective height of $\tilde{P_{t}}$ with respect to $\iota_{H}$. 

If  $H=\tilde{H}_{\Delta}$, then 

 $$H(\iota_{H}(\tilde{P_{t}})) \le (m+1)m^{24(m+1)}A^{2(m+1)}((d+1)H(\iota))^{m+2}B^{d(m+2)} .$$	

 
 If $H=H_{\operatorname{p}}$, then 

 $$H(\iota_{H}(\tilde{P_{t}})) \le  m^{24}(d+1)^{2}H(\iota)^{2}B^{2d} .$$	



\end{prop}

\begin{proof}
	\textbf{Case 1: $P_{t}$ lifts to $C_{\tilde{H}_{\Delta}}(K)$}
	
	
	As noted above, working on the quotient $C_{\tilde{H}_{\Delta}}$ we have $\iota_{\tilde{H}_{\Delta}}$ embeds $C_{\tilde{H}_{\Delta}}$ into $\P^{7}$ by composing $\iota^{\prime}_{\tilde{H}_{\Delta}}$ with the Segre embedding. A point $\tilde{P_{t}}=(E_{t},E_{t}^{\prime}, \psi: E_{t}[m] \stackrel{\sim}{\rightarrow} E_{t}^{\prime}[m])$ that is a lift of $P_{t}=(E_{t},E_{t}^{\prime})$ embedded into $\P^{7}$ as following:
	
	$$(E_{t},E_{t}^{\prime}, \psi) \to  F \times j(E_{t}) \times j(E_{t}^{\prime}) \hookrightarrow [Fj(E_{t})j(E_{t}^{\prime}), Fj(E_{t}), Fj(E_{t}^{\prime}),\cdots, 1].$$
	
	Since our ultimate goal is to count rational points parametrized by $t \in K$, we need a formula relating the height of $\iota_{H}(\tilde{P_{t}})$ with heights of $F$, $j(E_{t})$ and $j(E_{t}^{\prime})$. By Lemma \ref{ProductHeights}, we have
	\begin{equation}
	\label{3}
		H(\iota_{H}(\tilde{P_{t}}))=H(F)H(j(E_{t}))H(j(E_{t}^{\prime})).
	\end{equation}
	Let $i$ be an integer between $1$ and $m+1$ such that 
	$$ H(j(E_{t,i})j(E_{t,i}^{\prime}))= \max_{1 \le k \le m+1} H(j(E_{t,k})j(E_{t,k}^{\prime})).$$
	  By Definition \ref{F} and Corollary \ref{pazc} and Lemma \ref{5.3}, together with the fact that for any $\alpha, \beta, \alpha_{1} \cdots \alpha_{r} \in \overline{\Q}$,
	  $$H(\alpha\beta) \le H(\alpha)H(\beta)$$
	  and $$H(\alpha_{1}+\cdots\alpha_{r}) \le rH(\alpha_{1})\cdots H(\alpha_{r}),$$
	  we have 
	
	\begin{align}
	H(F) & =H(j(E_{t,1})j(E^{\prime}_{t,1})+ \cdots +j(E_{t, m+1})j(E^{\prime}_{t, m+1})) \\
	& \le (m+1) H(j(E_{t,i}))^{m+1}H(j(E_{t,i}^{\prime}))^{m+1} \\
	& \le (m+1)m^{24(m+1)}A^{2(m+1)}H(j(E_{t}))^{m+1}H(j(E_{t}^{\prime}))^{m+1}  \\
	& \le (m+1)m^{24(m+1)}A^{2(m+1)}((d+1)H(\iota))^{m+1}B^{d(m+1)} .
\end{align}
The constant $A$ comes from Corollary \ref{pazc}, and the first part of the lemma follows from $(\ref{3})$.

\vspace{1em}
\noindent \textbf{Case 2: $P_{t}$ lifts to one of the maximal parabolic quotient surfaces}	

Recall the definition of $H_{\operatorname{p}}$ (\ref{Hprod}) and $\iota_{H_{\operatorname{p}}}$ (\ref{iotaHp}). As in the previous case, Lemma \ref{ProductHeights} implies that 
 
 $$H(\iota_{H_{\operatorname{p}}}(\tilde{P_{t}}))=H(j(E_{t}))H(j(\tilde{E_{t}}))H(j(E_{t}^{\prime}))H(j(\tilde{E_{t}^{\prime}}))$$
 By Corollary \ref{pazc} and Lemma \ref{5.3}, where $\tilde{E}_{t}$ is $m$-isogenous to $E_{t}$ and $\tilde{E}^{\prime}_{t}$ is $m$-isogenous to $E_{t}^{\prime}$, we have 

$$H(\iota_{H_{\operatorname{p}}}(\tilde{P_{t}})) \ll m^{24}(d+1)^{2}H(\iota)^{2}B^{2d} .$$
 
		
\end{proof}

\end{document}
